\subsubsection*{Question 1}

$\dfrac{1}{3} \times \dfrac{6}{5}-\dfrac{1}{5} =  \dfrac{1 \times 2 \times 3}{3 \times 5} - \dfrac{1}{5} = \dfrac{2}{5} - \dfrac{1}{5} = \dfrac{1}{5}$

Bonne réponse : \textbf{C.~~} $\dfrac{1}{5}$


\subsubsection*{Question 2}

$\begin{aligned}
	3 x+2= -5 x+4
	&\iff 3x + 5x + 2 = 4 \\
	&\iff  8x + 2 =  4  \\
	&\iff  8x  =  2 \\
	&\iff x = \dfrac{2}{8} = \dfrac{1}{4}
\end{aligned}$

Bonne réponse : \textbf{C.~~} $x= \dfrac{1}{4}$

\subsubsection*{Question 3}

$(2x - 3)(2x + 3) = 2x\times 2x + 2x\times 3 + (-3)\times 2x + (-3)\times 3 = 4x^2 - 9$

Bonne réponse : \textbf{D.~~}$4x^2 - 9$

\subsubsection*{Question 4}

$5~\% = \dfrac{5}{100} = \dfrac{20}{400}$. L'effectif total des animaux est 400.

Bonne réponse : \textbf{B.~~} 400

\subsubsection*{Question 5}
$\np[km]{30} = \np[h]{0.5} \times \np[km/h]{60}$ et $\np[h]{0.5} = \np[min]{30}$.

Bonne réponse : \textbf{A.~~} 30

\subsubsection*{Question 6}
Augmenter de $5~\%$ par an c'est multiplier par \np{1.05} chaque année. Au bout de trois ans cela revient à multiplier par $\np{1.05}\times\np{1.05}\times\np{1.05} = \np{1.05}^3$.

Bonne réponse : \textbf{B.~~} $1,05^{3}$


\subsubsection*{Question 7}
Par lecture graphique, pour $x = 0$, $y = -1$. Seules les réponses A et B peuvent convenir.

Par lecture graphique, pour $x = 4$, $y = 0$. Par calcul : $4\times 4 - 1 = 15$  et $\dfrac{1}{4}\times 4 - 1 = 0$.

Bonne réponse : \textbf{B.~~} $y=\dfrac{1}{4} x - 1$


\subsubsection*{Question 8}
Bonne réponse : \textbf{A.~~}$f$ est croissante sur $[-2 ; 0]$
